\section{Reproducibility and Audit Trail}
\label{sec:reproducibility}

Scientific choices are stored as data rather than hidden in the report text.
The canonical schema, normalisation prompts, transition rules, lexical frames,
item-bank design, diagnostic acceptance rule, fold, candidate family,
selection obligations, predefined policies, oracle, and KT parameters are
independently versioned.  Executors validate these artifacts before applying
them.  Appendix~\ref{sec:artifact-map} maps each conceptual component to its
repository location.

Every run records the resolved experiment, random seed, external source
identity, Git commit and dirty state, and signatures of upstream inputs.  Stage
reuse is opt-in: a child experiment names a parent and starting stage, and the
runner refuses to copy earlier outputs if any recursively referenced input has
changed.  There is no implicit cache.  Model-assisted stages additionally
retain prompts, raw and parsed responses, invocation metadata, and errors for
each attempt.

The source contract declares a SHA-256 beginning
\texttt{e38c4f959f188150\ldots} and ending \texttt{c2488486cd}, and a total of
1,222 records; the full digest remains in the executable experiment file.
Because the snapshot is consult-only, it is located through
\texttt{GRAMMAR\_KT\_EGP\_SOURCE} or a declared data-root variable and is not
redistributed.  A digest or record-count mismatch is fatal.

The working model configuration currently requests
\texttt{gpt-5.6-sol} at medium reasoning effort in a fresh read-only context
for both normalisation and item diagnostics.  The service snapshot and all
decoding parameters are not pinned.  Exact external re-execution is therefore
not guaranteed even with identical local files.  Retaining every response and
repeating selected units makes this limitation inspectable; it does not remove
it.

The final verification boundary comprises four command-level checks:
\begin{itemize}
  \raggedright
  \item execute the complete \texttt{unittest} discovery suite with
        \texttt{PYTHONPATH=src};
  \item generate \texttt{base} with \texttt{scripts/run.py} and require a
        clean result from \texttt{scripts/validate.py};
  \item run each child condition explicitly from \texttt{kc\_selection}; and
  \item use \texttt{scripts/compare.py} at the compositional stage to verify
        fixed inputs and calculate paired contrasts.
\end{itemize}
Equivalent runs are required for the other conditions in
Table~\ref{tab:ontology-conditions}.  Inspection commands trace a descriptor,
cell, item, Q-matrix row, ordinary event, or probe through its saved evidence.
Automated tests verify implementation and leakage boundaries; the research
audit notebook separately records linguistic and methodological review.  A
green test suite does not override an unreviewed notebook or a failed run
validator.
